\section{Separation rates and the cost of outcome tails}

Holding the primitive smoothness tuple fixed gives a direct comparison of the precision attainable under the finite conditional moment restriction and the bounded-outcome benchmark. The comparison concerns the population heterogeneity distance and the capped critical radii defined in \cref{obj:def:critical-radius,obj:def:bounded-radius}. Two exponents determine the original-record separation rate: one depends on effect smoothness and moment order, while the other also reflects the smoothness of the propensity and baseline mean.

We collect the scales and their finite-sample multipliers before stating the comparison. Write \(C^{\mathrm{att}}_v\) and \(N^{\mathrm{att}}_v\) for the fully evaluable attainment multiplier and nonempty-power threshold specified in \cref{obj:thm:original-record-attainable-rate}; evaluation at \((2,w)\) gives their matched bounded counterparts. The following definition combines these attainment quantities with the lower-bound multipliers.

\begin{definitionv}{def:sharp-frontier-scales}[Separation scales \(\rho_n(v)\) and \(\rho_n^{\mathrm b}(w)\)]
The separation scales and lower multipliers are
\[
 \rho_n(v)=n^{-E(p)},\qquad
 \rho_n^{\mathrm b}(w)=n^{-E(2)},\qquad
 c_w^{\mathrm b}=c_{(2,w)},
\]
\[
 c_v=
 \begin{cases}
 c_0^{\mathrm e},&F(p)\ge1,\\
 \min\!\left\{k_v,\,
 c_0^{\mathrm e}N_{\mathrm{low}}(v)^{-(E_0(p)-E_4(p))}\right\},
 &F(p)<1.
 \end{cases}
\]
For an admissible public tuple \(v=(p,\alpha,\beta,\gamma)\in\mathcal V\) and its matched smoothness tuple \(w=(\alpha,\beta,\gamma)\), the quantities in these formulas are defined by
\[
 q=\frac{p-1}{p},\qquad S=\alpha+\beta,\qquad
 D_p=1+2\alpha+\frac{\beta}{q}+\frac{S}{2\gamma},
\]
\[
 E_0(p)=\frac{2\gamma q}{2\gamma+q},\qquad
 E_4(p)=\frac{2S}{D_p},\qquad
 F(p)=\frac{2\alpha}{p-1}+\frac{\beta}{q}+\frac{S}{2\gamma},
 \qquad E(p)=\min\{E_0(p),E_4(p)\},
\]
\[
 H_{\mathrm{fine}}=2^{40},\qquad
 N_{\mathrm{low}}(v)=
 \left\lceil\max\left\{2,32^{D_p\gamma/(2S)}\right\}\right\rceil,
\]
\[
 k_v=2^{-17}(2H_{\mathrm{fine}})^{-S},\qquad
 c_0^{\mathrm e}=\frac{4^{-\gamma}}{16\sqrt3}.
\]
The arguments \(p\) in \(E_0,E_4,F,E\) vary with \((\alpha,\beta,\gamma)\) fixed; \(c_{(2,w)}\) is \(c_v\) evaluated at \(v=(2,\alpha,\beta,\gamma)\).

The upper multipliers and sample-size thresholds are the fully evaluable attainment quantities:
\[
 C_v=C^{\mathrm{att}}_v,\qquad N_v=N^{\mathrm{att}}_v,\qquad
 C_w^{\mathrm b}=C^{\mathrm{att}}_{(2,w)},\qquad
 N_w^{\mathrm b}=N^{\mathrm{att}}_{(2,w)}.
\]
The test \(\phi^*_{n,v}\) is the total original-record test specified by the explicit score construction, and its bounded-model counterpart is
\[
 \phi^{*,\mathrm b}_{n,w}=\phi^*_{n,(2,w)}.
\]
\end{definitionv}

The original-model scale \leanref{sym:rho}{\ensuremath{\rho_n(v)}} and the matched bounded scale \leanref{sym:rhobound}{\ensuremath{\rho_n^{\mathrm b}(w)}} summarize the sample-size exponents. Their lower multipliers, \leanref{sym:clower}{\ensuremath{c_v}} and \leanref{sym:clowerbound}{\ensuremath{c_w^{\mathrm b}}}, and upper multipliers, \leanref{sym:Cupper}{\ensuremath{C_v}} and \leanref{sym:Cupperbound}{\ensuremath{C_w^{\mathrm b}}}, retain the finite-sample content of the comparison. The associated public thresholds \leanref{sym:Npublic}{\ensuremath{N_v}} and \leanref{sym:Npublicbound}{\ensuremath{N_w^{\mathrm b}}} place the upper separations below a fixed attainable heterogeneity distance. Thus the power statements below are attached to an explicit nonempty separation domain.

To measure the cost of outcome tails at the same smoothness tuple, we use the ratio of the capped radii.

\begin{definitionv}{def:comparison-handle}[Capped-radius ratio $\Delta_n(v)$]
\[
\Delta_n(v)=\frac{r_n^*(v)}{r_n^{*,\mathrm b}(w)},
\qquad w=(\alpha,\beta,\gamma).
\]
The numerator and denominator are the capped critical separation radii for the original and bounded original-record experiments, respectively, at the same sample size \(n\) and matched smoothness tuple \(w\).
\end{definitionv}

The ratio \leanref{sym:Delta}{\ensuremath{\Delta_n(v)}} in \cref{obj:def:comparison-handle} compares experiments with the same mean-effect target and error thresholds. We distinguish persistent enlargement of the original-model radius from a comparison that remains bounded over all sample sizes.

\begin{definitionv}{def:tail-region}[Tail region \(\mathcal T\)]
The tail region \(\mathcal T\) consists of admissible public tuples \(v=(p,\alpha,\beta,\gamma)\), with moment exponent \(p\) and primitive smoothness exponents \(\alpha,\beta,\gamma\), for which
\[
\mathcal T
=
\left\{
v\in\mathcal V:
p<2,\ 
\lim_{n\to\infty}\Delta_n(v)=\infty
\right\},
\]
where \(\mathcal V\) is the admissible tuple domain and \(\Delta_n(v)\) is the original-to-bounded capped-radius ratio at matched smoothness.
\end{definitionv}

\begin{definitionv}{def:equality-region}[Equality region \(\mathcal E\)]
The equality region \(\mathcal E\) is defined using \(\mathcal V\), the admissible domain of public moment and smoothness tuples, by
\[
\mathcal E
=
\left\{
v\in\mathcal V:
\sup_{n\ge 2}\Delta_n(v)<\infty
\right\},
\]
where \(\Delta_n(v)\) is the original-to-bounded capped-radius ratio at matched smoothness.
\end{definitionv}

The tail region \leanref{sym:TailRegion}{\ensuremath{\mathcal T}} and equality region \leanref{sym:EqualityRegion}{\ensuremath{\mathcal E}} describe the behavior of this matched comparison. The equality region concerns bounded radius ratios; its definition permits different finite-sample constants.

Before comparing rates, the bounded benchmark can be identified exactly with its signed-binary submodel. Let \(\lambda\) denote Lebesgue probability measure on \([0,1]\), and let \leanref{sym:dzero}{\ensuremath{d_0=1/(8\sqrt{2})}} be the fixed positive witness distance supplied by \cref{obj:lem:causal-nonempty}. The next result shows that both benchmarks have the same attainable heterogeneity distances and the same minimax testing risks.

\begin{propositionv}{prop:bounded-submodel-comparison}[Bounded and binary submodel comparison]
Let \(v=(p,\alpha,\beta,\gamma)\) be an admissible parameter tuple as in \cref{obj:def:model}, and let \(w=(\alpha,\beta,\gamma)\). Use the model classes in \cref{obj:def:binary-model,obj:def:bounded-model,obj:def:model}. Let \(\lambda\) denote Lebesgue probability measure on \([0,1]\), and define the heterogeneity distance, its model suprema, and the witness distance by
\[
d(P)=
\left[
\int_{[0,1]}
\left(\tau_P(x)-\int_{[0,1]}\tau_P(z)\,\lambda(dz)\right)^2
\,\lambda(dx)
\right]^{1/2},
\]
\[
D_w^{\mathrm{bin}}=\sup_{P\in\mathcal M^{\mathrm{bin}}_w}d(P),
\qquad
D_w^{\mathrm b}=\sup_{P\in\mathcal M^{\mathrm b}_w}d(P),
\qquad
D_v=\sup_{P\in\mathcal M_v}d(P),
\qquad
d_0=\frac{1}{8\sqrt{2}}.
\]
Then
\[
\mathcal M^{\mathrm{bin}}_w\subseteq\mathcal M^{\mathrm b}_w
\subseteq\mathcal M_v,
\qquad
D_w^{\mathrm{bin}}=D_w^{\mathrm b}=D_v\ge d_0.
\]
For every integer \(n\ge2\) and every real \(r\) satisfying \(0<r<D_v\), the testing risks of \cref{obj:def:binary-risk,obj:def:bounded-risk,obj:def:testing-risk} satisfy
\[
B_n^{\mathrm{bin}}(w,r)=B_n^{\mathrm b}(w,r)\le B_n(v,r).
\]
For every integer \(n\ge2\), the capped critical radii of \cref{obj:def:binary-radius,obj:def:bounded-radius,obj:def:critical-radius} satisfy
\[
0<r_n^{*,\mathrm{bin}}(w)
=r_n^{*,\mathrm b}(w)\le r_n^*(v).
\]

For every \(P\in\mathcal M_{(2,\alpha,\beta,\gamma)}\), its signed-binary conversion belongs to \(\mathcal M^{\mathrm{bin}}_w\). This conversion retains the distribution of \((X,A)\) and uses conditional signed-binary outcomes whose probabilities of \(+1\) in arms \(0\) and \(1\), respectively, are
\[
\frac{1+m_{0,P}(x)}{2},
\qquad
\frac{1+m_{1,P}(x)}{2},
\]
with the complementary probabilities assigned to \(-1\). It preserves the propensity \(e_P\), the baseline mean \(m_{0,P}\), the effect \(\tau_P\), and the heterogeneity distance \(d(P)\). The converted law satisfies the null constancy condition of \cref{obj:ass:null-constancy} if and only if the original law does.

For every \(n\in\mathbb N\), every test \(\phi\) in the sense of \cref{obj:def:testing-risk}, and every \(P\in\mathcal M^{\mathrm b}_w\), independently convert each original record by retaining \((X,A)\) and drawing an outcome sign with conditional probabilities
\[
\Pr(+1\mid X,A,Y)=\frac{1+Y}{2},
\qquad
\Pr(-1\mid X,A,Y)=\frac{1-Y}{2}.
\]
Integrating this conversion randomization into \(\phi\) gives a rejection-probability map on the original experiment. Its expected rejection probability under \(P\) equals the expected rejection probability of \(\phi\) under the signed-binary conversion of \(P\).
\end{propositionv}

The sign randomization in \cref{obj:prop:bounded-submodel-comparison} preserves the conditional outcome means, and hence the effect function, its heterogeneity distance, and null membership. Integrating the randomization into a test transfers its rejection probabilities to the bounded experiment. The resulting equality of minimax risks holds throughout the common nonempty separation domain \(0<r<D_v\), and equality of capped radii holds for every \(n\ge2\). This gives an exact decision-theoretic benchmark for the mean-effect target.

The lower-bound comparison uses finite mixtures of original-record laws. Their public tuning selects a copula-frame construction in the rough-confounding regime once its sample-size threshold is reached, and an opposite paired-tent construction otherwise. The following selection rule fixes the priors used in the main result; their support laws and complete-record comparison bounds are developed in \cref{obj:lem:paired-tent-full-record-lower,obj:def:copula-frame-priors,obj:lem:copula-frame-legality,obj:lem:full-record-copula-component-bound,obj:lem:rough-full-record-sharp-lower}.



The selected null and alternative priors \leanref{sym:piNull}{\ensuremath{\pi_{0,n,v}}} and \leanref{sym:piAlt}{\ensuremath{\pi_{1,n,v}}} in \cref{obj:def:converse-handle} induce the complete independent-record mixtures \leanref{sym:MixNull}{\ensuremath{\mathbb P_{0,n,v}}} and \leanref{sym:MixAlt}{\ensuremath{\mathbb P_{1,n,v}}}. Their rare-mark probability \leanref{sym:rarity}{\ensuremath{\varepsilon_{n,v}}} and outcome magnitude \leanref{sym:magnitude}{\ensuremath{L_{n,v}}} describe how the lower witnesses use the finite moment allowance.

For the next statement, \(\mathsf R_n(P,\phi)\) denotes expected rejection in the original-record experiment, \(Q_{a,P}(dy\mid x)\) denotes the conditional outcome probability kernel in arm \(a\), and \leanref{sym:TV}{\ensuremath{\operatorname{TV}}} denotes the supremum of the absolute difference in event probabilities. The threshold \(N(v)\) identifies the sample sizes at which every bounded-model finite prior pair at the original-model lower separation has mixture total variation at least one half. The theorem combines the matched separation bounds, the power guarantees, and this tail-essential comparison.



The exponent comparison in \cref{obj:thm:heavy-tail-comparison} explains the two statistical regimes. When \(F(p)\ge1\), the effect-smoothness exponent \(E_0(p)\) determines the separation scale: moment order and effect smoothness govern its sample-size dependence. When \(F(p)<1\), the rough-confounding exponent \(E_4(p)\) determines the scale, bringing the propensity and baseline smoothness exponents into the rate. The identity for \(E_4(p)-E_0(p)\) makes \(F(p)=1\) their exact equality boundary. Across both regimes, the strict inequality \(E(p)<E(2)\) for \(p<2\) translates the finite moment allowance into a polynomial enlargement of the matched critical radius.

The two denominators reflect different balances. At histogram rank \(J\), approximating the effect contributes the usual Hölder term of order \(J^{-\gamma}\); clipping at level \(T\) contributes a moment-tail term of order \(T^{1-p}\), while the clipped quadratic variability grows with \(T^{2-p}\). Balancing these terms gives \(E_0(p)\). When the nuisance functions are rough, the ordered-pair correction removes the first-order propensity error, leaving a product of the fine-scale propensity and baseline approximation errors. Its rank cost depends on \(2\alpha+\beta/q\): propensity regularity enters twice through treatment weighting and projection correction, whereas the baseline appears once and is amplified by the tail factor \(1/q\). The additional term \(S/(2\gamma)\) records the coarse resolution needed to retain the effect signal. Correction-specific clipping prevents the fine increments from paying the main cutoff at every resolution, which is what makes the multiresolution covariance sum compatible with \(E_4(p)\). These statements summarize the verified mean and covariance bounds in \cref{obj:lem:original-score-mean-ledger,obj:lem:original-score-covariance-ledger}; their full calculations are in the attainment appendix.

For example, \(v=(3/2,1,1,1)\) has \(F(p)=8\), so the effect channel governs with \(E=E_0=2/7\). In contrast, \(v=(3/2,1/20,1/20,1)\) has \(F(p)=2/5\), so rough confounding governs with \(E=E_4=2/13\). The first tuple uses smooth nuisance functions and the second deliberately makes them much rougher than the effect. These statistical regimes are distinct from the procedure's representation choice in \cref{obj:def:explicit-score-ledger}, which uses the conditions \(S\ge\gamma\) or \(\alpha\ge q\) to choose between a single correction and a multiresolution expansion.

The regimes can be summarized by substituting the formulas in \cref{obj:def:sharp-frontier-scales}. In the following table, \(q=(p-1)/p\) and \(S=\alpha+\beta\), with the primitive smoothness tuple held fixed.
\begin{center}
\begin{tabular}{lll}
\hline
Model and regime & Condition & Separation exponent \\
\hline
Original, effect smoothness
& \(F(p)>1\)
& \(\displaystyle E_0(p)=\frac{2\gamma q}{2\gamma+q}\) \\[6pt]
Original, equality boundary
& \(F(p)=1\)
& \(\displaystyle E_0(p)=E_4(p)=\frac{2\gamma q}{2\gamma+q}\) \\[6pt]
Original, rough confounding
& \(F(p)<1\)
& \(\displaystyle E_4(p)=\frac{2S}{1+2\alpha+\beta/q+S/(2\gamma)}\) \\[6pt]
Bounded, effect smoothness or boundary
& \(F(2)\ge1\)
& \(\displaystyle E_0(2)=\frac{2\gamma}{4\gamma+1}\) \\[6pt]
Bounded, rough confounding
& \(F(2)<1\)
& \(\displaystyle E_4(2)=\frac{2S}{1+2S+S/(2\gamma)}\) \\[6pt]
\hline
\end{tabular}
\end{center}
Here \(F(2)=2S+S/(2\gamma)\). Each row describes the power of \(n\) in the corresponding scale; the finite-sample radius bounds retain the multipliers and capping convention of \cref{obj:thm:heavy-tail-comparison}. The original and bounded exponents are selected separately at their respective moment orders.

For every fixed admissible tuple, the two-sided bounds on \(\Delta_n(v)\) quantify the cost of tails through \(n^{E(2)-E(p)}\). The moment-order-two face has equal original and bounded exponents and a uniformly bounded ratio over \(n\ge2\), together with the ordering \(\Delta_n(v)\ge1\). Every admissible tuple with \(p<2\) instead belongs to the tail region, where the ratio diverges. This distinction between moment orders is separate from the equality boundary \(F(p)=1\), which identifies the meeting of the two rate mechanisms within a given moment class.

The finite-sample cap remains part of the interpretation. The radius inequalities hold for every \(n\ge2\), while the power guarantees for the attaining tests \leanref{sym:phistar}{\ensuremath{\phi^*_{n,v}}} and \leanref{sym:phistarbound}{\ensuremath{\phi^{*,\mathrm b}_{n,w}}} from \cref{obj:def:sharp-frontier-scales} apply when their upper separations lie below the corresponding model distance suprema. The public thresholds ensure this condition through the fixed witness distance \(d_0\). Size is controlled over the whole null at every \(n\ge2\). On compact subsets of the admissible exponent domains, common lower and upper multipliers and a common bound on the public thresholds make these comparisons locally uniform.

Finally, the tail-essential witnesses identify the statistical content of the enlargement. For \(p<2\), their outcome magnitudes grow with sample size, while the selected null and separated alternative mixtures remain within total variation one half in the complete-record experiment. At the same original-model lower separation, every bounded-model finite prior pair has mixture total variation at least one half once the tuple-dependent threshold is reached. Thus the witnesses connect the slower finite-moment rate to outcome tails while preserving the stipulated primitive restrictions and the mean-effect target. The full support-law constructions and the quantified finite-prior comparison appear in the proof appendix through \cref{obj:def:marked-table,obj:def:frame-handle,obj:def:copula-frame-priors,obj:thm:tail-essential-full-record-converse}.
